%------------------------------------------------------------------------------%
\documentclass[fleqn,utf8,aspectratio=169,12pt]{beamer}

\usepackage{gaal}

\title{Operações com Matrizes}

%------------------------------------------------------------------------------%
\begin{document}

\addtitle

%------------------------------------------------------------------------------%
\section{Igualdade, Soma, Subtração e Produto por Escalar}

%------------------------------------------------------------------------------%
\begin{frame}{Igualdade de matrizes}
  \onslide<+->{Duas matrizes são iguais se}
  \begin{enumerate}[<+->]
    \item possuem a mesma ordem
    \item os elementos correspondentes são iguais
  \end{enumerate}

  \onslide<+->{%
  Assim,
  \[
    A = B
  \]}
  \onslide<+->{%
  se, e somente se,
  \[
    a_{ij} = b_{ij}
    \qquad\qquad
    \forall i = 1, \dots, m \quad\text{e}\quad j = 1, \dots, n
  \]}
\end{frame}

%------------------------------------------------------------------------------%
%------------------------------------------------------------------------------%
\begin{question}
Determine $x$ e $y$ sabendo que
\[
  \begin{pmatrix}
    2x & 3   \\
    4  & y-1
  \end{pmatrix}
  =
  \begin{pmatrix}
    8 & 3 \\
    4 & 5
  \end{pmatrix}
\]
\end{question}

%------------------------------------------------------------------------------%
\begin{resolution}{Solução}
  \onslide<+->{Comparando os elementos correspondentes}
  \begin{columns}
  \column{0.4\linewidth}
    \begin{align*}
      \onslide<+->{2x & = 8           \\}
      \onslide<+->{x  & = \frac{8}{2} \\}
      \onslide<+->{x  & = 4}
    \end{align*}
  \column{0.4\linewidth}
    \begin{align*}
      \onslide<+->{y-1 & = 5     \\}
      \onslide<+->{y   & = 5 + 1 \\}
      \onslide<+->{y   & = 4}
    \end{align*}
  \end{columns}
\end{resolution}

%------------------------------------------------------------------------------%

%------------------------------------------------------------------------------%
\begin{question}
  Determine $x$, $y$ e $z$ sabendo que
  \[
    \begin{pmatrix}
      x+1 & 2y  \\
      3   & z-2
    \end{pmatrix}
    =
    \begin{pmatrix}
      5 & 8 \\
      3 & 7
    \end{pmatrix}
  \]
\end{question}

%------------------------------------------------------------------------------%
\begin{resolution}{Solução}
  \onslide<+->{Comparando os elementos correspondentes}
  \begin{columns}
  \column{0.33\linewidth}
    \begin{align*}
      \onslide<+->{x+1 & = 5     \\}
      \onslide<+->{x   & = 5 - 1 \\}
      \onslide<+->{x   & = 4}
    \end{align*}
  \column{0.33\linewidth}
    \begin{align*}
      \onslide<+->{2y & = 8           \\}
      \onslide<+->{y  & = \frac{8}{2} \\}
      \onslide<+->{y  & = 4}
    \end{align*}
  \column{0.33\linewidth}
    \begin{align*}
      \onslide<+->{z-2 & = 7     \\}
      \onslide<+->{z   & = 7 + 2 \\}
      \onslide<+->{z   & = 9}
    \end{align*}
  \end{columns}
\end{resolution}

%------------------------------------------------------------------------------%


%------------------------------------------------------------------------------%
\begin{frame}{Adição de Matrizes}
A soma de duas matrizes é definida \emph{apenas} quando elas possuem
a \emph{mesma ordem}

\pause
Se
\[
  A = \big( a_{ij} \big)_{m\times n}
  \qquad\text{e}\qquad
  B = \big( b_{ij} \big)_{m\times n}
\]
\pause
então
\[
  A + B
  \pause
  = \big( a_{ij} + b_{ij} \big)_{m\times n}
\]
\pause
ou seja, somamos os elementos correspondentes
\end{frame}

%------------------------------------------------------------------------------%
%------------------------------------------------------------------------------%
\begin{question}
  Considere
  \quad
  \(
    A =
    \begin{pmatrix}
      1 & 2 \\
      3 & 4
    \end{pmatrix}
  \)
  \quad e \quad
  \(
    B =
    \begin{pmatrix}
      5 & -1 \\
      2 & 3
    \end{pmatrix}
  \)
  \quad
  avalie $A+B$
\end{question}

%------------------------------------------------------------------------------%
\begin{resolution}{Solução}
\begin{align*}
  \onslide<+->{A + B}
  \onslide<+->{& = \begin{pmatrix}
        1 & 2 \\
        3 & 4
      \end{pmatrix}
  +
      \begin{pmatrix}
        5 & -1 \\
        2 & 3
      \end{pmatrix}
  \\[2mm]}
  \onslide<+->{
  & \setlength{\arraycolsep}{12pt}
  = \begin{pmatrix}
    1+5 & 2+(-1) \\
    3+2 & 4+3
  \end{pmatrix}
  \\[2mm]}
  \onslide<+->{& =
  \begin{pmatrix}
    6 & 1 \\
    5 & 7
  \end{pmatrix}}
\end{align*}
\end{resolution}

%------------------------------------------------------------------------------%

%------------------------------------------------------------------------------%
\begin{question}
Calcule
\(
  A + B
\)
onde
\[
  A = \begin{pmatrix}
    2 & -3 & 1 \\
    4 & 0  & 5
  \end{pmatrix}
\qquad
  B = \begin{pmatrix}
    1  & 2 & -4 \\
    -2 & 3 & 1
  \end{pmatrix}
\]
\end{question}

%------------------------------------------------------------------------------%
\begin{resolution}{Solução}
\begin{align*}
  \onslide<+->{A + B}
  \onslide<+->{& = \begin{pmatrix}
    2 & -3 & 1 \\
    4 & 0  & 5
  \end{pmatrix}
  +
  \begin{pmatrix}
    1  & 2 & -4 \\
    -2 & 3 & 1
  \end{pmatrix}
  \\[2mm]}
  \onslide<+->{
  & \setlength{\arraycolsep}{12pt}
  = \begin{pmatrix}
    2 + 1    & -3 + 2 & 1 + (-4) \\
    4 + (-2) & 0 + 3  & 5 + 1
  \end{pmatrix}
  \\[2mm]}
  \onslide<+->{& = \begin{pmatrix}
  3 & -1 & -3\\
  2 & 3 & 6
  \end{pmatrix}}
\end{align*}

\end{resolution}

%------------------------------------------------------------------------------%


%------------------------------------------------------------------------------%
\begin{frame}{Subtração de Matrizes}
A subtração só é definida para matrizes de \emph{mesma ordem}

\pause
Se
\[
  A = \big( a_{ij} \big)_{m\times n}
  \qquad\text{e}\qquad
  B = \big( b_{ij} \big)_{m\times n}
\]
\pause
\[
  A - B = \big( a_{ij} - b_{ij} \big)_{m\times n}
\]

\pause
Podemos escrever
\[
  A - B = A + (-B)
\]
\end{frame}

%------------------------------------------------------------------------------%
%------------------------------------------------------------------------------%
\begin{question}
Avalie a subtração de $A$ por $B$, sendo
\;
\(
  A =
  \begin{pmatrix}
    7 & 3  \\
    2 & -1
  \end{pmatrix}
\)
\; e \;
\(
  B=
  \begin{pmatrix}
    4  & 5 \\
    -2 & 3
  \end{pmatrix}
\)

\end{question}

%------------------------------------------------------------------------------%
\begin{resolution}{Solução}
\begin{align*}
  \onslide<+->{A - B}
  \onslide<+->{& = \begin{pmatrix}
    7 & 3  \\
    2 & -1
  \end{pmatrix}
  - \begin{pmatrix}
    4  & 5 \\
    -2 & 3
  \end{pmatrix}
  \\[2mm]}
  \onslide<+->{& = \begin{pmatrix}
    7-4    & 3-5  \\
    2-(-2) & -1-3
  \end{pmatrix}
  \\[2mm]}
  \onslide<+->{& = \begin{pmatrix}
    3 & -2 \\
    4 & -4
  \end{pmatrix}}
\end{align*}
\end{resolution}

%------------------------------------------------------------------------------%

%------------------------------------------------------------------------------%
\begin{question}
Calcule
\(
 A - B
\),
onde
\[
  A =
  \begin{pmatrix}
    5 & -2 & 4  \\
    1 & 3  & -1
  \end{pmatrix}
\qquad \text{ e } \qquad
  B =
  \begin{pmatrix}
    2 & 1  & -3 \\
    4 & -2 & 5
  \end{pmatrix}
\]
\end{question}

%------------------------------------------------------------------------------%
\begin{resolution}{Solução}
\begin{align*}
  \onslide<+->{A - B}
  \onslide<+->{& =
    \begin{pmatrix}
      5 & -2 & 4  \\
      1 & 3  & -1
    \end{pmatrix}
  - \begin{pmatrix}
      2 & 1  & -3 \\
      4 & -2 & 5
    \end{pmatrix}
  \\[2mm]}
  \onslide<+->{& =
    \begin{pmatrix}
      5 - 2 & -2 -1 & 4 -(-3) \\
      1 - 4 & 3 - (-2) & -1 -5
    \end{pmatrix}
  \\[2mm]}
  \onslide<+->{& = \begin{pmatrix}
    3  & -3 & 7  \\
    -3 & 5  & -6
  \end{pmatrix}}
\end{align*}
\end{resolution}

%------------------------------------------------------------------------------%


%------------------------------------------------------------------------------%
\begin{frame}{Multiplicação por Escalar}
  Se $c$ é um escalar (número real) e $A$ é uma matriz,
  \pause
  então
  \[
    cA
  \]
  \pause
  é obtida multiplicando \emph{cada elemento} de $A$ por $c$

  \pause
  Se
  \[
    A = \big( a_{ij}\big)
  \]
  \pause
  então
  \[
    cA = \big( ca_{ij} \big)
  \]
\end{frame}

%------------------------------------------------------------------------------%
%------------------------------------------------------------------------------%
\begin{question}
Considerando
\(
  A =
  \begin{pmatrix}
    1 & -2 \\
    3 & 4
  \end{pmatrix}
\)
avalie $3A$
\end{question}

%------------------------------------------------------------------------------%
\begin{resolution}{Solução}
\begin{align*}
  \onslide<+->{3A}
  \onslide<+->{& = 3 \begin{pmatrix}
    1 & -2 \\
    3 & 4
  \end{pmatrix}
  \\[2mm]}
  \onslide<+->{
  & \setlength{\arraycolsep}{12pt}
  = \begin{pmatrix}
    3 \times 1 & 3(-2) \\
    3 \times 3 & 3 \times 4
  \end{pmatrix}
  \\[2mm]}
  \onslide<+->{& = \begin{pmatrix}
    3 & -6 \\
    9 & 12
  \end{pmatrix}}
\end{align*}
\end{resolution}

%------------------------------------------------------------------------------%
%------------------------------------------------------------------------------%
\begin{question}
  Calcule
  \(
    -2A
  \), onde
  \(
    A =
    \begin{pmatrix}
      3 & -1 & 2  \\
      0 & 4  & -5
    \end{pmatrix}
  \)
\end{question}

%------------------------------------------------------------------------------%
\begin{resolution}{Solução}
\begin{align*}
  \onslide<+->{-2A}
  \onslide<+->{& = -2 \begin{pmatrix}
    3 & -1 & 2  \\
    0 & 4  & -5
  \end{pmatrix}
  \\[2mm]}
  \onslide<+->{
  & \setlength{\arraycolsep}{12pt}
  = \begin{pmatrix}
    -2\times 3 & (-2)(-1) & -2\times 2  \\
    -2\times 0 & -2\times 4  & (-2)(-5)
  \end{pmatrix}
  \\[2mm]}
  \onslide<+->{& = \begin{pmatrix}
    -6 & 2  & -4 \\
    0  & -8 & 10
  \end{pmatrix}}
\end{align*}
\end{resolution}

%------------------------------------------------------------------------------%


%------------------------------------------------------------------------------%
\begin{frame}{Propriedades}
Sejam $A$ e $B$ matrizes de ordem $m \times n$ e $\alpha$, $\beta \in \R$
\pause
\begin{enumerate}[<+->][(i)\;]
  \setlength\itemsep{1.3ex}
  \item $A + B = B + A$                               \tabto{67mm} comutatividade
  \item $(A + B) + C = A + (B + C)$                   \tabto{67mm} associatividade
  \item $A + \mathcal{O} = A$                         \tabto{67mm} elemento neutro na adição
  \item $\alpha (A + B) = \alpha A + \alpha B$        \tabto{67mm} distributividade do escalar
  \item $(\alpha + \beta)A = \alpha A + \beta A$      \tabto{67mm} distributividade da matriz
  \item $0A = \mathcal{O}$
  \item $(\alpha \beta)A = \alpha (\beta A)$
\end{enumerate}
\end{frame}

%------------------------------------------------------------------------------%
%------------------------------------------------------------------------------%
\begin{question}
Calcule
\(
  2A-3B
\)
para
\(
  A=
  \begin{pmatrix}
    1 & 2  \\
    0 & -1
  \end{pmatrix}
\)
e
\(
  B=
  \begin{pmatrix}
    2 & -1 \\
    3 & 1
  \end{pmatrix}
\)
\end{question}

%------------------------------------------------------------------------------%
\begin{resolution}{Solução}
\[
  2A
  \pause
  \;=\;
  2 \begin{pmatrix}
    1 & 2  \\
    0 & -1
  \end{pmatrix}
  \pause
  \;=\;
  {\setlength{\arraycolsep}{12pt}
  \begin{pmatrix}
    2 \times 1 & 2 \times 2  \\
    2 \times 0 & 2(-1)
  \end{pmatrix}}
  \pause
  \;=\;
  \begin{pmatrix}
    2 & 4\\
    0 & -2
  \end{pmatrix}
\]
\pause
\[
  3B
  \pause\;=\;
  3 \begin{pmatrix}
    2 & -1 \\
    3 & 1
  \end{pmatrix}
  \pause
  \;=\;
  {\setlength{\arraycolsep}{12pt}
  \begin{pmatrix}
    3 \times 2 & 3(-1) \\
    3 \times 3 & 3 \times 1
  \end{pmatrix}}
  \pause
  \;=\;
  \begin{pmatrix}
    6 & -3 \\
    9 & 3
  \end{pmatrix}
\]

\pause
\smallskip
\[
  2A-3B
  \pause\;=\;
  \begin{pmatrix}
    2 & 4\\
    0 & -2
  \end{pmatrix}
  +
  \begin{pmatrix}
    6 & -3 \\
    9 & 3
  \end{pmatrix}
  \pause
  \;=\;
  {\setlength{\arraycolsep}{12pt}
  \begin{pmatrix}
    2 - 6 & 4 - (-3)\\
    0 - 9 & -2 - 3
  \end{pmatrix}}
  \pause
  \;=\;
  \begin{pmatrix}
    -4 & 7  \\
    -9 & -5
  \end{pmatrix}
\]
\end{resolution}

%------------------------------------------------------------------------------%

%------------------------------------------------------------------------------%
\begin{question}
Considerando
\(
  A =
  \begin{pmatrix}
    2 & -1 \\
    3 & 0
  \end{pmatrix}
\)
\; e \;
\(
  B =
  \begin{pmatrix}
    1  & 2 \\
    -1 & 4
  \end{pmatrix}
\), \;
calcule
\(
  3A-2B
\)
\end{question}

%------------------------------------------------------------------------------%
\begin{resolution}{Solução}
\[
  3A
  \;=\;
  3 \begin{pmatrix}
    2 & -1 \\
    3 & 0
  \end{pmatrix}
  \;=\;
  \begin{pmatrix}
    3 \times 2 & 3(-1) \\
    3 \times 3 & 3 \times 0
  \end{pmatrix}
  \;=\;
  \begin{pmatrix}
    6 & -3 \\
    9 & 0
  \end{pmatrix}
\]

\[
  2B
  \;=\;
  2 \begin{pmatrix}
    1  & 2 \\
    -1 & 4
  \end{pmatrix}
  \;=\;
  \begin{pmatrix}
    2 \times 1  & 2 \times 2 \\
    2(-1) & 2 \times 4
  \end{pmatrix}
  \;=\;
  \begin{pmatrix}
    2  & 4 \\
    -2 & 8
  \end{pmatrix}
\]

\[
  3A-2B
  \;=\;
  \begin{pmatrix}
    6 & -3 \\
    9 & 0
  \end{pmatrix}
  -
  \begin{pmatrix}
    2  & 4 \\
    -2 & 8
  \end{pmatrix}
    \;=\;
  \begin{pmatrix}
    6-2    & -3-4 \\
    9-(-2) & 0-8
  \end{pmatrix}
  \;=\;
  \begin{pmatrix}
    4  & -7 \\
    11 & -8
  \end{pmatrix}
\]

\end{resolution}

%------------------------------------------------------------------------------%


%------------------------------------------------------------------------------%
\section{Multiplicação}

%------------------------------------------------------------------------------%
\begin{frame}{Produto de Matrizes}
Se
\(
  A_{m\times \emph{k}}
\)
\; e \;
\(
  B_{\emph{k}\times n}
\)
\pause
então o produto
\[
  AB
\]
\pause
é definido e resulta em uma matriz
\[
  C_{m\times n} = AB
\]

\pause
\medskip
\emph{Condição}
\[
\text{número de colunas de }A
=
\text{número de linhas de }B
\]
\end{frame}

%------------------------------------------------------------------------------%
\begin{frame}{Definição}
  Dadas duas matrizes $A_{m \times \emph{k}}$ e $B_{\emph{k} \times n}$,
  \pause
  podemos definir a \emph{matriz produto}
  \[
    C_{m \times n} = AB
  \]
  \pause
  cujos elementos são obtidos pela fórmula
  \[
    c_{ij}
    \pause
    \;=\; a_{i1} b_{1j} + a_{i2} b_{2j} + \cdots + a_{ik} b_{kj}
    \pause
    \;=\; \sum_{l=1}^{k} a_{il} b_{lj}
  \]
  \pause
  para $i = 1, \dots, m$ e $j = 1, \dots, n$
\end{frame}

%------------------------------------------------------------------------------%
\begin{frame}{Matrizes de Ordem $2\times 2$}
Considere
\[
  A =
  \begin{pmatrix}
    a & b \\
    c & d
  \end{pmatrix}
  \qquad
  B =
  \begin{pmatrix}
    e & f \\
    g & h
  \end{pmatrix}
\]
\pause
Então
\[
  \setlength{\arraycolsep}{12pt}
  AB =
  \begin{pmatrix}
    ae+bg & af+bh \\
    ce+dg & cf+dh
  \end{pmatrix}
\]
\end{frame}

%------------------------------------------------------------------------------%
%------------------------------------------------------------------------------%
\begin{question}
Considerando
\(
  A =
  \begin{pmatrix}
    1 & 2 \\
    3 & 4
  \end{pmatrix}
\)
\(
  B =
  \begin{pmatrix}
    5 & 6 \\
    7 & 8
  \end{pmatrix}
\),
calcule $AB$
\end{question}

%------------------------------------------------------------------------------%
\begin{resolution}{Solução}
\begin{align*}
  \onslide<+->{AB}
  \onslide<+->{& =
  \begin{pmatrix}
    1 & 2 \\
    3 & 4
  \end{pmatrix}
  \begin{pmatrix}
    5 & 6 \\
    7 & 8
  \end{pmatrix}
  \\[2mm]}
  \onslide<+->{& =
  \begin{pmatrix}
    1\times 5 + 2\times 7 & 1\times 6 + 2\times 8 \\
    3\times 5 + 4\times 7 & 3\times 6 + 4\times 8
  \end{pmatrix}
  \\[2mm]}
  \onslide<+->{& =
  \begin{pmatrix}
    19 & 22 \\
    43 & 50
  \end{pmatrix}}
\end{align*}
\end{resolution}

%------------------------------------------------------------------------------%

%------------------------------------------------------------------------------%
\begin{question}
Calcule $BA$ com as matrizes anteriores
\end{question}

%------------------------------------------------------------------------------%
\begin{resolution}{Solução}
\begin{align*}
  \onslide<+->{BA}
  \onslide<+->{& =
  \begin{pmatrix}
    5 & 6 \\
    7 & 8
  \end{pmatrix}
  \begin{pmatrix}
    1 & 2 \\
    3 & 4
  \end{pmatrix}
  \\[2mm]}
  \onslide<+->{& =
  \begin{pmatrix}
    5 \times 1 + 6 \times 3 & 5 \times 2 + 6 \times 4 \\
    7 \times 1 + 8 \times 3 & 7 \times 2 + 8 \times 4
  \end{pmatrix}
  \\[2mm]}
  \onslide<+->{& =
  \begin{pmatrix}
    23 & 34 \\
    31 & 46
  \end{pmatrix}}
  \onslide<+->{
  \\[2mm]
  & \neq
  \begin{pmatrix}
  19 & 22\\
  43 & 50
  \end{pmatrix}
  = AB}
\end{align*}
\end{resolution}

%------------------------------------------------------------------------------%

%------------------------------------------------------------------------------%
\begin{question}
Calcule $AB$, onde
\(
  A=
  \begin{pmatrix}
    2 & -1 & 3 \\
    0 & 4  & 1
  \end{pmatrix}
\)
e
\(
  B=
  \begin{pmatrix}
    1  & 2 \\
    3  & 0 \\
    -1 & 4
  \end{pmatrix}
\)
\end{question}

%------------------------------------------------------------------------------%
\begin{resolution}{Solução}
Observe que a matriz
\[
  M = A_{2\times3}B_{3\times2}
\]
é definida e produz uma matriz $2\times2$
\end{resolution}

%------------------------------------------------------------------------------%
\begin{resolution}{Solução}
\begin{align*}
  \onslide<+->{M = AB}
  \onslide<+->{& =
  \begin{pmatrix}
    2 & -1 & 3 \\
    0 & 4  & 1
  \end{pmatrix}
  \begin{pmatrix}
    1  & 2 \\
    3  & 0 \\
    -1 & 4
  \end{pmatrix}
  \\[2mm]}
  \onslide<+->{
  & \setlength{\arraycolsep}{12pt}
  = \begin{pmatrix}
    2 \times 1 +       (-1)3 + 3(-1) & 2\times 2 +     (-1)0 + 3 \times 4 \\
    0 \times 1 +  4 \times 3 + 1(-1) & 0\times 2 + 4\times 0 + 1 \times 4
  \end{pmatrix}
  \\[2mm]}
  \onslide<+->{& =
  \begin{pmatrix}
  -4 & 16\\
  11 & 4
  \end{pmatrix}}
\end{align*}
\end{resolution}

%------------------------------------------------------------------------------%

%------------------------------------------------------------------------------%
\begin{question}
Determine a transposta de
\(
  A =
  \begin{pmatrix}
    1 & 2 & 3 \\
    4 & 5 & 6
  \end{pmatrix}
\)
\end{question}

%------------------------------------------------------------------------------%
\begin{resolution}{Solução}
  \[
    A^t
    \pause
    \;=\;
    \begin{pmatrix}
      1 & 2 & 3 \\
      4 & 5 & 6
    \end{pmatrix}^t
    \pause
    \;=\;
    \begin{pmatrix}
      1 & 4 \\
      2 & 5 \\
      3 & 6
    \end{pmatrix}
  \]
\end{resolution}

%------------------------------------------------------------------------------%


%------------------------------------------------------------------------------%
\begin{frame}{Propriedades do Produto Matricial}

\begin{enumerate}[<+->][(i)\;]
\itemsep=1em
\item $A(BC) = (AB)C$
      \tabto{85mm} associatividade
\item \(I_m A_{m \times n} = A_{m \times n}\)
      \quad e \quad
      \(A_{m \times n} I_n = A_{m \times n}\)
      \tabto{85mm} elemento neutro
\item $A(B + C) = AB + AC$
\item $(B + C)A = BA + CA$
\item $\alpha(AB) = (\alpha A)B = A(\alpha B)$
      \tabto{85mm} para qualquer escalar $\alpha$
\end{enumerate}

\bigskip
\onslide<+->{\emph{Atenção}: em geral \quad
\(
  AB \neq BA
\)}
\end{frame}

%------------------------------------------------------------------------------%
\begin{frame}{Produto de Hadamard ou Produto Elemento a Elemento}
Para matrizes de mesma ordem,
\[
  A = \big( a_{ij} \big)_{m\times n}
  \qquad\text{e}\qquad
  B = \big( b_{ij} \big)_{m\times n}
\]
\pause
o produto de Hadamard é
\[
  A\circ B = \big( a_{ij} b_{ij} \big)
\]

\pause
Este produto é implementado em algumas linguagens de programação

\pause
\emph{Não utilizaremos} este produto nesta disciplina
\end{frame}

%------------------------------------------------------------------------------%
\section{Transposição}

%------------------------------------------------------------------------------%
\begin{frame}{Matriz Transposta}
  A transposta de uma matriz $A$ é obtida trocando suas
  linhas por colunas

  \pause
  A transposta de $A$ é indicada por
  \[
    A^T
    \qquad\text{ ou }\qquad
    A^t
  \]

  \pause
  Se
  \[
    A = \big( a_{ij} \big)
  \]
  \pause
  então
  \[
    A^T = A^t = \big( a_{ji} \big)
  \]
\end{frame}

%------------------------------------------------------------------------------%
\begin{frame}{Propriedades da Transposição}
  \begin{enumerate}[<+->][(i)\;]
  \itemsep=1.3em
    \item $\left(A^t\right)^t = A$
    \item $(A + B)^{t} = A^{t} + B^{t}$
    \item $(AB)^{t} = B^t A^t$
    \item $(\alpha A)^{t} = \alpha A^{t} $
  \end{enumerate}
\end{frame}

%------------------------------------------------------------------------------%
%------------------------------------------------------------------------------%
\begin{question}
Determine a matriz transposta de
\(
  A =
  \begin{pmatrix}
    2  & -1 & 0 \\
    3  & 4  & 5 \\
    -2 & 1  & 6
  \end{pmatrix}
\)
\end{question}

%------------------------------------------------------------------------------%
\begin{resolution}{Solução}
  \[
    A^t
    \pause
    \;=\;
    \begin{pmatrix}
      2  & -1 & 0 \\
      3  & 4  & 5 \\
      -2 & 1  & 6
    \end{pmatrix}^t
    \pause
    \;=\;
    \begin{pmatrix}
      2  & 3 & -2 \\
      -1 & 4 & 1  \\
      0  & 5 & 6
    \end{pmatrix}
  \]
\end{resolution}

%------------------------------------------------------------------------------%

%------------------------------------------------------------------------------%
\begin{question}
  Considerando
  \(
    A =
    \begin{pmatrix}
      2 & 1 \\
      1 & 3
    \end{pmatrix}
  \)
  e
  \(
    x =
    \begin{pmatrix}
      2  \\
      -1
    \end{pmatrix}
  \),
  calcule
  \(
    v = x^t A x
  \)
\end{question}

%------------------------------------------------------------------------------%
\begin{resolution}{Solução}
  \[
    x^t
    \pause
    \;=\;
    \begin{pmatrix}
      2  \\
      -1
    \end{pmatrix}^t
    \pause
    \;=\;
    \begin{pmatrix}
      2 & -1
    \end{pmatrix}
  \]

  \pause
  \[
    Ax
    \pause
    \;=\;
    \begin{pmatrix}
      2 & 1 \\
      1 & 3
    \end{pmatrix}
    \begin{pmatrix}
      2  \\
      -1
    \end{pmatrix}
    \pause
    \;=\;
    \begin{pmatrix}
      2 \times 2 + 1(-1) \\
      1 \times 2 + 3(-1)
    \end{pmatrix}
    \pause
    \;=\;
    \begin{pmatrix}
      3  \\
      -1
    \end{pmatrix}
  \]

  \pause
  \[
    v
    \;=\;
    x^tAx
    \pause
    \;=\;
    \begin{pmatrix}
      2 & -1
    \end{pmatrix}
    \begin{pmatrix}
      3  \\
      -1
    \end{pmatrix}
    \pause
    \;=\;
    2 \times 3 + (-1)(-1)
    \pause
    \;=\;
    7
  \]
\end{resolution}

%------------------------------------------------------------------------------%

\input{ex/B-04-Operacoes-ex-17}
%------------------------------------------------------------------------------%
\begin{question}
Considerando os vetores \;
\(
  u =
  \begin{pmatrix}
    2  \\
    -1 \\
    3
  \end{pmatrix}
\) \;
e
\; \(
  v =
  \begin{pmatrix}
    1  \\
    4  \\
    -2
  \end{pmatrix}
\),\;
calcule
\(
  uv^t
\)
\end{question}

%------------------------------------------------------------------------------%
\begin{resolution}{Solução}
\begin{align*}
  \onslide<+->{uv^t}
  \onslide<+->{& =
  \begin{pmatrix}
    2  \\
    -1 \\
    3
  \end{pmatrix}
  \begin{pmatrix}
    1  \\
    4  \\
    -2
  \end{pmatrix}^t}
  \onslide<+->{\;=\;
  \begin{pmatrix}
    2  \\
    -1 \\
    3
  \end{pmatrix}
  \begin{pmatrix}
    1 & 4 & -2
  \end{pmatrix}
  \\}
  \onslide<+->{& =
  \begin{pmatrix}
    2\times 1 & 2\times 4 & 2\times (-2) \\
    (-1)1     & (-1)4     & (-1)(-2)     \\
    3\times 1 & 3\times 4 & 3\times (-2)
  \end{pmatrix}
  \\}
  \onslide<+->{& =
  \begin{pmatrix}
    2  & 8  & -4 \\
    -1 & -4 & 2  \\
    3  & 12 & -6
  \end{pmatrix}}
\end{align*}
\end{resolution}

%------------------------------------------------------------------------------%


%------------------------------------------------------------------------------%
\begin{frame}{Matrizes Simétricas}
  Uma matriz quadrada, $A_n$, é dita \emph{simétrica},
  \pause
  se
  \[
    A = A^t
  \]
  \pause
  isto é, se
  \quad
  \(
    a_{ij} = a_{ji}
    \qquad
    \forall \quad 1 \leq i,j \leq n
  \)

  \pause
  \bigskip
  Ela á \emph{antissimétrica},
  \pause
  se
  \[
    A = -A^t
  \]
  \pause
  isto é, se
  \quad
  \(
    a_{ij} = -a_{ji}
    \qquad
    \forall \quad 1 \leq i,j \leq n
  \)
\end{frame}
%------------------------------------------------------------------------------%
%------------------------------------------------------------------------------%
\begin{question}
Determine os valores de $x$, $y$ e $z$ para que a matriz
\[
  A =
  \begin{pmatrix}
    2 & x  & 4 \\
    3 & 5  & y \\
    z & -1 & 7
  \end{pmatrix}
\]
seja simétrica
\end{question}

%------------------------------------------------------------------------------%
\begin{resolution}{Solução}
\begin{align*}
  \onslide<+->{A & = A^t \\}
  \onslide<+->{
  \begin{pmatrix}
    2 & x  & 4 \\
    3 & 5  & y \\
    z & -1 & 7
  \end{pmatrix}
  & =
  \begin{pmatrix}
    2 & 3 & z  \\
    x & 5 & -1 \\
    4 & y & 7
  \end{pmatrix}}
\end{align*}

\bigskip
\[
  \onslide<+->{x =  3 \qquad}
  \onslide<+->{y = -1 \qquad}
  \onslide<+->{z = 4}
\]
\end{resolution}

%------------------------------------------------------------------------------%

%------------------------------------------------------------------------------%
\begin{question}
Determine os valores de $x$, $y$, $z$ e $w$ para que a matriz
\[
A=
  \begin{pmatrix}
    0  & x & 4 \\
    -3 & 0 & y \\
    z  & 2 & w
  \end{pmatrix}
\]
seja antissimétrica
\end{question}

%------------------------------------------------------------------------------%
\begin{resolution}{Solução}
\begin{columns}
\column{0.4\linewidth}
\begin{align*}
  \onslide<+->{A & = -A^t}
  \\
  \onslide<+->{
  \begin{pmatrix}
    0  & x & 4 \\
    -3 & 0 & y \\
    z  & 2 & w
  \end{pmatrix}
  & = -
  \begin{pmatrix}
    0 & -3 &  z \\
    x &  0 &  2 \\
    4 &  y &  w
  \end{pmatrix}}
  \\
  \onslide<+->{
  \begin{pmatrix}
    0  & x & 4 \\
    -3 & 0 & y \\
    z  & 2 & w
  \end{pmatrix}
  & =
  \begin{pmatrix}
     0 &  3 & -z \\
    -x &  0 & -2 \\
    -4 & -y & -w
  \end{pmatrix}}
\end{align*}
\column{0.4\linewidth}
\begin{align*}
  \onslide<+->{x & = 3 \\}
  \onslide<+->{y & = -2 \\}
  \onslide<+->{z & = -4 \\[7mm]}
  \onslide<+->{w & = -w \\}
  \onslide<+->{w & = 0}
\end{align*}
\end{columns}
\end{resolution}

%------------------------------------------------------------------------------%


%------------------------------------------------------------------------------%
\section{Lista Mínima}

%------------------------------------------------------------------------------%
\begin{frame}{Lista Mínima}
  \begin{enumerate}
    \item
  \end{enumerate}

  \vfill
  Atenção: A prova é baseada no livro, não nas apresentações
\end{frame}

%------------------------------------------------------------------------------%
\end{document}
%------------------------------------------------------------------------------%
